\begin{lexlean}{Guide}
\useglossary{lexlean.std.bool@1.2.0}
\useglossary{lexlean.std.nat@1.2.0}
\useglossary{prism.arch@1.0.0}
\importmodule{Foundation.View.Text.V3.Model}
\title{UORust workbench view}
\begin{semanticmodule}
definition{
  name: purpose,
  parameters: [],
  result: String,
  body: "UORust reads a Cargo workspace and writes a complete Prism model of it. This page runs UORust's route table, the verified function that decides which Prism record each kind of Rust fact becomes, or refuses the fact."
}

definition{name: runningText, parameters: [], result: String, body: "Routing the fact."}

definition{
  name: routedText,
  parameters: [],
  result: String,
  body: "Routed: every kind names this route. The lines give the route, the fact kind, the Prism record kind and its facet."
}

definition{
  name: errorText,
  parameters: [],
  result: String,
  body: "Not routed: the fact is refused or the request is invalid. The details name the reason; its fix is listed under refusals and errors."
}

definition{
  name: overflowText,
  parameters: [],
  result: String,
  body: "The answer is too large for the page. Route this fact from the command line."
}

definition{
  name: sections,
  parameters: [],
  result: (List Foundation.View.Text.V3.Model::Section),
  body: (cons (record Foundation.View.Text.V3.Model::Section {id := "arch", title := "Architecture facet", paragraphs := (cons "crate-bin-target: UR-R01 to Foundation.Arch::Component. A [[bin]] target is a component of kind releaseBinary." (cons "crate-lib-target: UR-R02 to Foundation.Arch::Component. A [lib] target is a component of kind service." (cons "crate-doc: UR-R03 to Foundation.Arch::Concern. A crate-level doc comment states the concern the crate answers." (cons "package-author: UR-R04 to Foundation.Arch::Stakeholder. An authors entry or CODEOWNERS owner is a stakeholder." (cons "dependency-edge: UR-R05 to Foundation.Arch::Edge. A workspace-internal dependency is a control-flow edge between components." (nil String))))))}) (cons (record Foundation.View.Text.V3.Model::Section {id := "features", title := "Features and journeys", paragraphs := (cons "clap-subcommand: UR-R06 to Foundation.Arch::Feature. A clap subcommand is a feature realized by the call it dispatches to." (cons "pub-fn: UR-R07 to Feature.calls. A public function is an interface call a feature names." (cons "pub-trait: UR-R08 to Feature.calls. A public trait method is an interface call a feature names." (cons "example-target: UR-R09 to Foundation.Arch::Journey. An [[example]] target is a journey through the features it calls." (cons "test-fn: UR-R10 to Feature.acceptance. A #[test] function is an acceptance row of the feature it exercises." (cons "doctest: UR-R11 to Feature.acceptance. A doc test is an acceptance row of the item it documents." (nil String)))))))}) (cons (record Foundation.View.Text.V3.Model::Section {id := "behavior", title := "Behavior facet", paragraphs := (cons "error-enum: UR-R12 to LexLean.inductive. An enum that implements Error is a typed failure, an inductive with one constructor per variant." (cons "fn-body: UR-R13 to LexLean.definition. A function body is lifted to a LexLean definition; what does not lift is residue in frontier.toml." (nil String)))}) (cons (record Foundation.View.Text.V3.Model::Section {id := "design", title := "Design views", paragraphs := (cons "state-enum: UR-R26 to Foundation.Arch::UmlState. A fieldless enum matched in transitions is a UML state machine; each variant a state." (cons "call-edge: UR-R27 to Foundation.Arch::UmlMessage. A call across crates is a synchronous message of the UML sequence view." (cons "module: UR-R28 to Foundation.Arch::OpmThing. A module is an OPM object; its public functions are the processes that zoom into it." (nil String))))}) (cons (record Foundation.View.Text.V3.Model::Section {id := "sec", title := "Security facet", paragraphs := (cons "unsafe-block: UR-R14 to Foundation.Sec::Threat. An unsafe block is a memory-safety threat (elevation of privilege)." (cons "ffi-extern: UR-R15 to Foundation.Sec::Threat. An extern block crosses the language boundary: a threat (elevation of privilege)." (cons "fs-use: UR-R16 to Foundation.Sec::Asset. A std::fs use is a file-system capability: an asset." (cons "net-use: UR-R17 to Foundation.Sec::Asset. A std::net use is a network capability: an asset." (cons "process-use: UR-R18 to Foundation.Sec::Asset. A std::process use is a process capability: an asset." (cons "env-use: UR-R19 to Foundation.Sec::Asset. A std::env use reads the environment: an asset." (cons "forbid-unsafe: UR-R20 to Foundation.Sec::ApplicationSecurityControl. #![forbid(unsafe_code)] is a control whose evidence is the compiler gate." (cons "advisory: UR-R21 to Foundation.Sec::Risk. A RustSec advisory against a locked package is a risk on the supply chain." (nil String)))))))))}) (cons (record Foundation.View.Text.V3.Model::Section {id := "qual", title := "Quality facet", paragraphs := (cons "lint-config: UR-R23 to Foundation.Qual::QualityRequirement. A [lints] table entry is a maintainability requirement." (cons "rust-version: UR-R24 to Foundation.Qual::QualityRequirement. The rust-version (MSRV) field is a portability requirement." (cons "bench-target: UR-R25 to Foundation.Qual::QualityMeasure. A [[bench]] target is a performance-efficiency measure." (nil String))))}) (cons (record Foundation.View.Text.V3.Model::Section {id := "release", title := "Release facet", paragraphs := (cons "lock-package: UR-R22 to System.SpdxPackage. A Cargo.lock package is an SPDX package of the release's bill of materials." (nil String))}) (nil Foundation.View.Text.V3.Model::Section))))))))
}

definition{
  name: errors,
  parameters: [],
  result: (List Foundation.View.Text.V3.Model::ErrorHint),
  body: (cons (record Foundation.View.Text.V3.Model::ErrorHint {code := "unknown-fact", title := "unknown-fact", explanation := "No route names this kind of Rust fact, so UORust records it as an open claim instead of guessing a record.", recovery := "Use one of the kinds in the route table, or propose a new route in model/routes.toml.", fields := (nil String)}) (cons (record Foundation.View.Text.V3.Model::ErrorHint {code := "ambiguous", title := "ambiguous", explanation := "Two kinds of the same fact name different routes, so the fact could become two different records.", recovery := "Classify the fact by one kind, or by kinds that all name the same route.", fields := (nil String)}) (cons (record Foundation.View.Text.V3.Model::ErrorHint {code := "blank-fact", title := "blank-fact", explanation := "A line of the request is empty.", recovery := "Write one kind per line with no empty line and no trailing line feed.", fields := (nil String)}) (cons (record Foundation.View.Text.V3.Model::ErrorHint {code := "empty", title := "empty", explanation := "The request is empty.", recovery := "Enter at least one kind of Rust fact, such as pub-fn.", fields := (nil String)}) (cons (record Foundation.View.Text.V3.Model::ErrorHint {code := "too-large", title := "too-large", explanation := "The request is longer than 512 bytes.", recovery := "Route one fact at a time; a fact has only a few kinds.", fields := (nil String)}) (cons (record Foundation.View.Text.V3.Model::ErrorHint {code := "invalid-utf8", title := "invalid-utf8", explanation := "The request is not valid UTF-8 text.", recovery := "Enter the kinds as plain UTF-8 text.", fields := (nil String)}) (nil Foundation.View.Text.V3.Model::ErrorHint)))))))
}

definition{
  name: examples,
  parameters: [],
  result: (List Foundation.View.Text.V3.Model::Example),
  body: (cons (record Foundation.View.Text.V3.Model::Example {id := "error-enum", label := "Error enum", description := "An enum that implements Error becomes a typed failure in the behavior facet.", values := (cons "error-enum" (nil String)), expectedResponse := "{\"version\":1,\"status\":\"ok\",\"operation\":\"route\",\"message\":\"Routed: every kind names this route. The lines give the route, the fact kind, the Prism record kind and its facet.\",\"rows\":4,\"columns\":1,\"cells\":[\"route UR-R12\",\"fact error-enum\",\"record LexLean.inductive\",\"facet behavior\"],\"details\":[],\"request\":\"error-enum\",\"response\":\"route UR-R12\\nfact error-enum\\nrecord LexLean.inductive\\nfacet behavior\",\"error\":\"\"}", explanation := "route UR-R12"}) (cons (record Foundation.View.Text.V3.Model::Example {id := "unsafe-block", label := "Unsafe block", description := "An unsafe block becomes a security threat.", values := (cons "unsafe-block" (nil String)), expectedResponse := "{\"version\":1,\"status\":\"ok\",\"operation\":\"route\",\"message\":\"Routed: every kind names this route. The lines give the route, the fact kind, the Prism record kind and its facet.\",\"rows\":4,\"columns\":1,\"cells\":[\"route UR-R14\",\"fact unsafe-block\",\"record Foundation.Sec::Threat\",\"facet sec\"],\"details\":[],\"request\":\"unsafe-block\",\"response\":\"route UR-R14\\nfact unsafe-block\\nrecord Foundation.Sec::Threat\\nfacet sec\",\"error\":\"\"}", explanation := "route UR-R14"}) (cons (record Foundation.View.Text.V3.Model::Example {id := "same-route", label := "Two kinds, one route", description := "Kinds that name the same route agree.", values := (cons "test-fn\ntest-fn" (nil String)), expectedResponse := "{\"version\":1,\"status\":\"ok\",\"operation\":\"route\",\"message\":\"Routed: every kind names this route. The lines give the route, the fact kind, the Prism record kind and its facet.\",\"rows\":4,\"columns\":1,\"cells\":[\"route UR-R10\",\"fact test-fn\",\"record Feature.acceptance\",\"facet features\"],\"details\":[],\"request\":\"test-fn\\ntest-fn\",\"response\":\"route UR-R10\\nfact test-fn\\nrecord Feature.acceptance\\nfacet features\",\"error\":\"\"}", explanation := "route UR-R10"}) (cons (record Foundation.View.Text.V3.Model::Example {id := "ambiguous", label := "Ambiguous fact", description := "An enum that is both an error and a state machine is refused, not defaulted.", values := (cons "error-enum\nstate-enum" (nil String)), expectedResponse := "{\"version\":1,\"status\":\"error\",\"operation\":\"route\",\"message\":\"Not routed: the fact is refused or the request is invalid. The details name the reason; its fix is listed under refusals and errors.\",\"rows\":0,\"columns\":0,\"cells\":[],\"details\":[\"refused ambiguous\",\"fact error-enum\",\"fact state-enum\"],\"request\":\"error-enum\\nstate-enum\",\"response\":\"refused ambiguous\\nfact error-enum\\nfact state-enum\",\"error\":\"ambiguous\"}", explanation := "refused ambiguous"}) (cons (record Foundation.View.Text.V3.Model::Example {id := "unknown", label := "Unknown kind", description := "A kind no route names is refused as an open claim.", values := (cons "struct" (nil String)), expectedResponse := "{\"version\":1,\"status\":\"error\",\"operation\":\"route\",\"message\":\"Not routed: the fact is refused or the request is invalid. The details name the reason; its fix is listed under refusals and errors.\",\"rows\":0,\"columns\":0,\"cells\":[],\"details\":[\"refused unknown-fact\",\"fact struct\"],\"request\":\"struct\",\"response\":\"refused unknown-fact\\nfact struct\",\"error\":\"unknown-fact\"}", explanation := "refused unknown-fact"}) (nil Foundation.View.Text.V3.Model::Example))))))
}

definition{
  name: view,
  parameters: [],
  result: Foundation.View.Text.V3.Model::WorkbenchView,
  body: (record Foundation.View.Text.V3.Model::WorkbenchView {title := "UORust", heading := "UORust route table", purpose := (purpose), audience := "Rust developers who model a Cargo workspace in Prism, and Prism modelers who check what UORust emits.", guidance := (cons "Enter the kinds of one Rust fact, one per line, such as error-enum." (cons "Route it and read the first line: route, or the reason it is refused." (cons "Every kind in the route table below routes to exactly one Prism record kind." (cons "Two kinds that name different routes are refused as ambiguous, never resolved by a default." (nil String))))), initialOperation := "route", initialExample := "error-enum", operations := (cons (record Foundation.View.Text.V3.Model::Operation {id := "route", label := "Route a fact", summary := "Route the kinds of one Rust fact to the Prism record kind UORust emits for it.", help := "One kind per line; every kind must name the same route.", requestPrefix := "ui-route", fields := (cons (record Foundation.View.Text.V3.Model::Field {id := "facts", label := "Rust fact kinds", kind := (ctor Foundation.View.Text.V3.Model::FieldKind.Multiline), help := "One kind per line, such as pub-fn or unsafe-block; at most 503 bytes on this page.", placeholder := "pub-fn", initialValue := "error-enum", required := true, choices := (nil Foundation.View.Text.V3.Model::Choice)}) (nil Foundation.View.Text.V3.Model::Field)), examples := (examples)}) (nil Foundation.View.Text.V3.Model::Operation)), sections := (sections), links := (cons (record Foundation.View.Text.V3.Model::Link {label := "Source", httpsUrl := "https://github.com/afflom/uorust"}) (cons (record Foundation.View.Text.V3.Model::Link {label := "Route table", httpsUrl := "https://github.com/afflom/uorust/blob/main/model/routes.toml"}) (cons (record Foundation.View.Text.V3.Model::Link {label := "PrismPM", httpsUrl := "https://github.com/UOR-Foundation/prismpm"}) (nil Foundation.View.Text.V3.Model::Link)))), errors := (errors), statuses := (record Foundation.View.Text.V3.Model::StatusCopy {scriptRequired := "Enable JavaScript to route a fact.", initializing := "Loading the route table.", ready := "Ready.", running := (runningText), cancelling := "Cancelling.", cancelled := "Routing was cancelled.", timeout := "Routing took too long.", loadError := "The route table could not be loaded.", computeError := "Routing failed.", inputError := "Enter UTF-8 text of at most 503 bytes.", responseError := "The answer could not be displayed."}), design := (record Foundation.View.Text.V3.Model::Design {palette := (ctor Foundation.View.Text.V3.Model::Palette.Indigo), density := (ctor Foundation.View.Text.V3.Model::Density.Comfortable), width := (ctor Foundation.View.Text.V3.Model::Width.Balanced), colorScheme := (ctor Foundation.View.Text.V3.Model::ColorScheme.System)}), raw := (record Foundation.View.Text.V3.Model::RawEditor {requestPrefix := "ui-raw", label := "Raw request", help := "Send the exact request bytes the command line reads: one kind per line.", placeholder := "pub-fn", initialValue := "pub-fn", submitLabel := "Route request", responseLabel := "Answer", expanded := false}), response := (record Foundation.View.Text.V3.Model::DisplaySpec {title := "Answer", emptyText := "No answer yet.", rawLabel := "Exact answer", maximumCells := (64 : UInt32), maximumDetails := (64 : UInt32)}), execution := (record Foundation.View.Text.V3.Model::Execution {timeoutMilliseconds := (30000 : UInt32), cancelLabel := "Cancel", retryLabel := "Try again", wasmCachePolicy := (ctor Foundation.View.Text.V3.Model::WasmCachePolicy.NoStore)}), labels := (record Foundation.View.Text.V3.Model::WorkbenchLabels {navigation := "Sections", operation := "Operation", examples := "Examples", loadExample := "Load example", compute := "Route", guide := "Route table", links := "Links", fields := "Fields", expectedResponse := "Expected answer", details := "Details", resultTable := "Answer lines", row := "Line", column := "Column", request := "Request", reset := "Reset", skipToWorkbench := "Skip to the route table"})})
}
\end{semanticmodule}
\end{lexlean}
